\documentclass[12pt]{article}
\usepackage{amsmath, amssymb, amsthm}
\usepackage[english]{babel}
\usepackage{geometry}
\usepackage{mathrsfs}
\usepackage{physics}

\numberwithin{equation}{section}

% Other packages...
\newtheorem{theorem}{Theorem}[section]
\newtheorem{corollary}{Corollary}[theorem]
\newtheorem{lemma}[theorem]{Lemma}

\title{\textbf{Exercise Solutions of Chp 3 N. Goldenfeld Phase T. and RN}}
\author{Dr. Jed Pixley (solver) and Elisha Shmalo (writer)}
\date{\today}

\geometry{left=1in, right=1in, top=1in, bottom=1in}
\renewcommand{\baselinestretch}{1.25}

\begin{document}

\maketitle

\section{Chapter 3}

\subsection{Exercise 3.1}

Transfer matrix in $d=1$ Ising model with \textit{periodic} boundry conditions:

(a) We construct the matrix $S$ which diagonalises the transfer matrix $T$ such that 

\begin{equation}
    T' = S^{-1}TS
\end{equation}

is diagonal: Recall $h := \beta H, \ K := \beta J$ and 

\begin{equation}
    T = \left(\begin{matrix}
        e^{h+K} & e^{-K} \\
        e^{-K} & e^{-h + K}
    \end{matrix}\right).
\end{equation}

Define now the e-vals with the diagonal matrix 

\begin{equation}
    T' = \left(\begin{matrix}
        \lambda_+ &  \\
         & \lambda_-
    \end{matrix}\right).
\end{equation}

and it can easily be seen that 

\begin{equation}
    \lambda^2 - 2\lambda e^K\cosh(h) + 2\sinh(2K) = 0,
\end{equation}

where $\lambda \sim \lambda_{(\pm )}$. It follows then that

\begin{equation}
    \lambda = e^K [\cosh(h) \pm \sqrt{\sinh^2(h) + e^{-4K}} \ ].
\end{equation}

To diagonalize the matrix we must find e-vectors which involves solving 

\begin{equation}
    T \vec{v} = \lambda_{(\pm)}\vec{v}.
\end{equation}

Let $\vec{v}_{\pm} = (a_{\pm}, b_{\pm})$. Then, if we define 

\begin{equation}
    \cot(2\phi) = e^{2K}\sinh(h), 
\end{equation}

and plug away at a lot of algebra, we get to 

\begin{equation}
    \begin{aligned}
        b_+ &= a_+ \tan(\phi) \\ 
        a_- &= -b_- \cot(\phi)
    \end{aligned}
\end{equation}

if we require our e-vectors to be normalized this gives us 

\begin{equation}
    \begin{aligned}
        \vec{v}_+ &= (\cos(\phi), \sin(\phi)), \\
        \vec{v}_- &= (\sin(\phi), -\cos(\phi)),
    \end{aligned}
\end{equation}

so 

\begin{equation}
    S = \left(\begin{matrix}
        \cos(\phi) & \sin(\phi) \\
        \sin(\phi) & -\cos(\phi),
    \end{matrix}\right)
\end{equation}

where 

\begin{equation}
    \cot(2\phi) = e^{2K}\sinh(h).
\end{equation}

(b)

Note that 

\begin{equation}
    S^{-1} = S
\end{equation}

So that 

\begin{equation}
    S\sigma_zS^{-1} = \left(\begin{matrix}
        \cos^2(\phi) - \sin^2(\phi) & 2\cos(\phi)\sin(\phi) \\
        2\cos(\phi)\sin(\phi) & -(\cos^2(\phi) - \sin^2(\phi))
    \end{matrix}\right) = \left(\begin{matrix}
        \cos(2\phi) & \sin(2\phi) \\
        \sin(2\phi) & -\cos(2\phi)
    \end{matrix}\right)
\end{equation}

Thus 

\begin{equation}
    S\sigma_zS^{-1}(T')^N = \left(\begin{matrix}
        \cos(2\phi) & \sin(2\phi) \\
        \sin(2\phi) & -\cos(2\phi)
    \end{matrix}\right) \left(\begin{matrix}\lambda_1^N & \\ & \lambda_2^N\end{matrix}\right) = \left(\begin{matrix}
        \lambda_1^N\cos(2\phi) & \lambda_2^N\sin(2\phi) \\
        \lambda_1^N\sin(2\phi) & -\lambda_2^N\cos(2\phi)
    \end{matrix}\right) 
\end{equation}

Recall also that 

\begin{equation}
    \expval{S_i} = \frac{\text{Tr}(S^{-1\sigma_zS(T')^N})}{Z_n},
\end{equation}

where $Z_N = \lambda_1^N + \lambda_2^N$. THus we have 

\begin{equation}
    \expval{S_i} = \frac{\lambda_1^N \cos(2\phi) - \lambda_2\cos(2\phi)}{\lambda_1 + \lambda_2} = \frac{\cos(2\phi) - (\lambda_2/\lambda_1)^N\cos(2\phi)}{1 +  (\lambda_2/\lambda_1)^N}.
\end{equation}

Taking $N \to \infty$ this becomes 

\begin{equation}
    \expval{S_i} = \cos(2\phi).
\end{equation}

As shown in the book 

\begin{equation}
    \expval{S_i S_{i+j}} 1 /Z \text{Tr}(S^{-1}\sigma_zS(T')^j S^{-1}\sigma_z S(T')^{N-j}).
\end{equation}

Define 

\begin{equation}
    S^{-1}\sigma_z S = \left(\begin{matrix}
        e & g \\ f & k
    \end{matrix}\right)
\end{equation}

THen as $N \to \infty$

\begin{equation}
    \expval{S_i S_{i+j}} = e^2 + gf(\frac{\lambda_2}{\lambda_1})^j = \cos^2(2\phi) + \sin^2(2\phi)(\frac{\lambda_2}{\lambda_1})^j
\end{equation}

showing that the system is translationally invarient and 

\begin{equation}
    G(i, i+j) = \expval{S_iS_{i+j}} - \expval{S_i}\expval{S_{j}} = \sin^2(2\phi)(\frac{\lambda_2}{\lambda_1})^j
\end{equation}

(c) Recall 

\begin{equation}
    M = \sinh(h) / \sqrt{\sinh^2(h) + e^{-4k}}
\end{equation}

or alternativly 

\begin{equation}
    M(H) = \sinh(\beta H ) / \sqrt{\sinh^2(\beta H) + e^{-4k}}
\end{equation}


\end{document}